\section{Conclusions}
\label{sec:discussion}


This paper studied binary polynomial optimization (BPO) through the lens of its multilinear polytope and developed a decision-diagram (DD) framework for generating cutting planes over local structures. We presented a compact DD encoding of the multilinear set that considers only vertex variables and implicitly recovers the hyperedge variables along states and arcs. This encoding yields an extended network-flow formulation of $MP_G$ and supports a target-cut separation model adapted to the implicitly assigned hyperedge variables. We analyzed the size of this construction via a valid-antichain characterization of its width, which allowed us to identify hypergraph structures---including flowers and cycles---that admit bounded-width DDs, as well as a new family of hypergraphs for which BPO is solvable in polynomial time. Building on this machinery, we proposed three strategies for choosing the local support fed to the separator: a fixed-support baseline that mirrors known inequalities in the literature, an expanded-support rule that grows a seed by overlap, and a partition-support scheme based on section hypergraphs. Finally, we introduced a novel DD-based dimension and facet-audit procedure to certify the polyhedral quality of the generated cuts.

Computationally, on the LABS and VISION datasets, the fixed-support DD separator reproduces the behavior of the corresponding literature cuts, confirming that the target-cut model recovers comparable strength when given the same local information. The main gains, however, come from strengthening the support before separation: expanded and, especially, partition supports close substantially more of the root gap with far fewer cuts---on LABS the partition family closes up to $92.4\%$ of the root gap, and on VISION it closes the gap almost completely on every tested instance---and these stronger relaxations translate into a faster branch-and-bound, with the partition setting solving all VISION instances and giving the best practical performance on LABS. The facet audit corroborates this picture: strengthened supports yield cuts that are far more often maximal and facet-defining than those from the literature. Taken together, the results support the central message of the paper: a DD is best viewed not as another hand-derived family of inequalities, but as a black-box strengthening device that turns any chosen support into a facet-producing separation oracle.

Several limitations remain. The current implementation generates cuts only at the root, uses offline support choices, and relies on heuristic selection of structure. Future work should integrate the separator into a full branch-and-cut routine, combine DD cuts with the solver's native cuts, and develop adaptive criteria for selecting supports during the solve. Another promising direction is to use the relatively cheap $I_K$-facet audit as a cut-selection tool, prioritizing cuts that are not only violated but also locally facet-defining. Finally, the randomized nature of the PT family suggests studying multi-seed or adaptive partitioning schemes that preserve PT0's strong performance while improving robustness.

More broadly, the same approach could be applied to other binary optimization problems where strong inequalities arise from small but meaningful supports and, thus, enhance modern solvers such as Gurobi. Moreover, our procedure can be seen as a research tool in its own right. Specifically, it can help researchers identify interesting problem structures (i.e., for BPO or other problems), obtain strong cuts from them, and then infer properties and develop algorithms to generate those same cuts without needing to create a DD. 

